\clearpage
\section{Offline CI: Acceptance Bound to a Candidate}
\label{app:offline-ci-evidence}

The governance addendum in Section~\ref{app:milestone-20260719}
describes role and route coverage. A release gate also needs a way to
decide whether the evidence it has actually applies to the candidate
being evaluated. For the X Studio X expansion, this section proposes
one local gate for that purpose. Its input is a frozen local snapshot;
the result is a technical eligibility decision. The model assumes
trustworthy snapshot acquisition and makes no empirical claim about
the current implementation or any prior hackathon build.

\begin{definition}[Candidate and local evidence snapshot]
\label{def:offline-ci-snapshot}
Let $c=(s,a,p)$, where $s$ is the exact source manifest, $a$ is the
exact artifact manifest, and $p$ is the exact policy manifest, with
equality componentwise. Each manifest is a canonical finite byte
string representing its full policy-defined contents. The policy
determines which toolchain, dependency, build-input, and environment
identities are material and where they are bound. This is an exact
identity model; a digest-based implementation needs a separately
justified binding assumption.

A validated policy determines a finite set $Q_p$ of required check
identities and authorized producer and reviewer identities for each
check. An identity of a check binds its procedure and configuration
under $p$. Inputs are invalid if candidate or policy validation fails,
if the snapshot cannot be read and authenticated as a whole, or if
$Q_p$ is empty. An authenticated snapshot may contain no records.

Let $E$ be the finite collection of records in the frozen snapshot.
A receipt names a candidate, check, producer, and result (pass or
fail). A review names the exact receipt and includes an explicit
endorsement decision. A receipt is eligible for $(c,q)$ only when
it is well formed, authenticated, exactly bound to $(c,q)$, and
issued by an authorized producer for $q$. A review qualifies only
when it is well formed, authenticated, issued by an authorized
reviewer distinct from the receipt producer, and endorses the exact
eligible receipt. An invalid individual record supplies no evidence;
an invalid snapshot blocks the gate. Identity authenticity and the
trustworthiness of attestation verification are model hypotheses.
\end{definition}

\begin{theorem}[Exact-candidate acceptance and stale-evidence isolation]
\label{thm:offline-ci-acceptance}
For validated inputs in Definition~\ref{def:offline-ci-snapshot},
write $P_q$ for the existence of an eligible pass receipt with a
qualifying review, and $F_q$ for the existence of an eligible fail
receipt. A fail receipt needs no second review to veto readiness.
For invalid inputs the decision is \textsc{Blocked}. For validated
inputs, the gate applies the following rules in order:
\begin{enumerate}
\item \textsc{Blocked} if some $q\in Q_p$ has both $P_q$ and $F_q$.
\item \textsc{Rejected} if some $q\in Q_p$ has $F_q$.
\item \textsc{Blocked} if some $q\in Q_p$ lacks $P_q$.
\item \textsc{Ready} otherwise.
\end{enumerate}
The decision is single-valued, independent of record order and
duplicate copies, and \textsc{Ready} is equivalent to
\begin{equation}
\label{eq:offline-ci-ready}
   \forall q\in Q_p,\quad P_q\land\neg F_q .
\end{equation}
Adding records bound exclusively to other candidates leaves the
decision unchanged on an otherwise unchanged, valid snapshot.
In particular, when moving to a different candidate, a snapshot
containing only evidence for the previous candidate cannot make the
new candidate ready.
\end{theorem}

\begin{proof}
\begin{enumerate}
\item \textbf{Validity and domains.}
For a validated candidate, the components $s$, $a$, and $p$
are byte strings in the specified manifest domains.
The equality test applies to these three objects,
rather than to their names or their timestamps.
The policy determines $Q_p$ and the applicable authorization rules
before the evidence is considered.
In particular, the appearance of a check in a receipt cannot
add that check to $Q_p$, and its absence cannot remove a check.
Validation supplies a finite nonempty set $Q_p$
and a finite authenticated collection $E$.
If any validation premise fails, the outer rule returns
\textsc{Blocked} before forming any predicates.
Thus a partial or unauthenticated read is distinguished from
a valid snapshot with no evidence.
This distinction matters even if both inputs give the same decision:
only the latter can be used to report which checks lack evidence.

\item \textbf{Eligible witnesses.}
Fix an arbitrary $q\in Q_p$.
Let $A_q$ be the set of receipt contents from $E$
which are eligible for $(c,q)$ and report pass.
Let $B_q$ be the set of receipt contents from $E$
which are eligible for $(c,q)$ and report fail.
For $r\in A_q$, let $V_r$ be the set of contents of qualifying
reviews of $r$ in $E$.
These are finite sets by taking subsets of the finite snapshot
and identifying exact duplicate contents.
Define
\begin{align}
   A_q^+ &= \{r\in A_q : V_r\ne\varnothing\},
      \label{eq:offline-ci-reviewed}\\
   P_q &\iff A_q^+\ne\varnothing,
      \label{eq:offline-ci-pass}\\
   F_q &\iff B_q\ne\varnothing .
      \label{eq:offline-ci-fail}
\end{align}
Each predicate is Boolean because set membership is fixed
by the validated snapshot and its fixed verification rules.
A signed endorsement of a different receipt is outside $V_r$.
An endorsement from the producer of $r$ is outside $V_r$,
including when that producer is also authorized to review
another contributor's work.
If $A_q$ contains a receipt with no qualifying reviewer,
that receipt is absent from $A_q^+$.
It therefore cannot witness $P_q$.
Conversely, a witness in $A_q^+$ supplies both a pass receipt
and a qualifying review, by conjunction elimination on the
membership condition in \eqref{eq:offline-ci-reviewed}.
A failure only requires membership in $B_q$.
Thus omitting or delaying a failure review cannot remove its veto.

\item \textbf{Exhaustive per-check partition.}
For the fixed $q$, excluded middle gives either $P_q$
or $\neg P_q$, and either $F_q$ or $\neg F_q$.
Taking the two splits together gives four cases:
\begin{enumerate}
\item $(P_q,F_q)=(0,0)$.
The sets $A_q^+$ and $B_q$ are both empty.
There may be an eligible pass without a qualifying review,
or there may be no eligible receipts at all.
Either way, there is no witness for acceptance at $q$.
\item $(P_q,F_q)=(1,0)$.
There exists a reviewed pass and there is no eligible fail.
The existential test permits additional eligible passes,
including passes still awaiting review.
This case establishes acceptance for this check only.
\item $(P_q,F_q)=(0,1)$.
At least one eligible fail exists.
An unreviewed pass does not change this case.
The gate cannot use the lack of review to treat a known failure
as missing evidence.
\item $(P_q,F_q)=(1,1)$.
Two result classes are present for the same candidate and check.
The witnesses can come from the same producer on different
executions or from different authorized producers.
The gate presents the resulting conflict for resolution.
It does not infer a preferred result from order or frequency.
\end{enumerate}
The cases are pairwise disjoint because unequal Boolean pairs
differ in at least one coordinate.
They are jointly exhaustive by the two splits just established.
A difference in result between two candidates does not create
this conflict because the eligible sets are constructed at fixed $c$.
A difference between distinct checks of the same candidate
also does not create this conflict because each check has its own
eligible sets.

\item \textbf{Partition at the gate.}
For valid inputs introduce the sets
\begin{align}
   C &= \{q\in Q_p : P_q \land F_q\},
      \label{eq:offline-ci-conflict}\\
   N &= \{q\in Q_p : F_q\},
      \label{eq:offline-ci-negative}\\
   M &= \{q\in Q_p : \neg P_q\}.
      \label{eq:offline-ci-missing}
\end{align}
These are subsets of $Q_p$ by set comprehension.
In particular $C\subseteq N$ by conjunction elimination.
The sets can overlap. We therefore partition by the following
ordered tests, each conjoined with the negation of earlier tests:
\begin{enumerate}
\item $C\ne\varnothing$.
A conflict witness exists, so the decision is \textsc{Blocked}.
This also covers a simultaneous fail-only result at another
check, and a simultaneous missing result at a third check.
\item $C=\varnothing$ and $N\ne\varnothing$.
There is an eligible failure with no conflicting reviewed
pass at any check, and the decision is \textsc{Rejected}.
This holds even if $M\ne\varnothing$.
\item $C=N=\varnothing$ and $M\ne\varnothing$.
There is neither conflict nor eligible failure,
but at least one required check lacks a reviewed pass.
The decision is \textsc{Blocked}.
\item $C=N=M=\varnothing$.
Every check has a reviewed pass and none has an eligible fail.
The decision is \textsc{Ready}.
\end{enumerate}
For disjointness, two different cases disagree on the first
set whose emptiness distinguishes them.
For exhaustiveness, test emptiness of $C$, then of $N$ if $C$
is empty, then of $M$ if both previous sets are empty.
Each test has exactly two outcomes since the corresponding
set is fixed.
The partition with the initial validation branch added
therefore assigns exactly one decision to every input.
The two occurrences of \textsc{Blocked} for valid inputs
represent different reasons and can retain their separate
witness sets in a decision record.

\item \textbf{Necessity and sufficiency.}
Suppose the decision is \textsc{Ready}.
Then the valid-input partition put it in its fourth case.
In particular $M=\varnothing$ and $N=\varnothing$.
For arbitrary $q\in Q_p$, $M=\varnothing$ implies
$q\notin M$ and therefore $P_q$ by the definition of $M$.
Likewise, $q\notin N$ implies $\neg F_q$.
Universally generalizing this conjunction over $Q_p$
yields \eqref{eq:offline-ci-ready}.
Conversely, assume \eqref{eq:offline-ci-ready}
for a validated input.
If $M$ were nonempty then a member would witness
$\neg P_q$, contrary to \eqref{eq:offline-ci-ready}.
If $N$ were nonempty then a member would witness
$F_q$, also a contradiction.
Hence $M=N=\varnothing$.
Because $C\subseteq N$, $C$ is empty as well.
The fourth case gives \textsc{Ready}.
The same formula evaluated on an empty $Q_p$ would be
vacuously true, which explains why the validation condition
must precede the equivalence.

\item \textbf{Order and multiplicity.}
Consider two valid representations of a frozen snapshot
having the same set of record contents.
For every $q\in Q_p$ and receipt content $r$,
membership of $r$ in $A_q$ and $B_q$ depends only
on whether $r$ appears and on its fixed eligibility properties.
For a review, presence and qualification are likewise
properties of its contents and the exact receipt it names.
Therefore $V_r$ is equal in the two representations.
It follows by set extensionality that $A_q^+$ is equal,
then by the two existence tests that $P_q$ and $F_q$
are equal in the two representations.
This argument holds for every $q$.
Consequently $C$, $N$, and $M$ are equal.
The gate gives the same decision.
A permutation and the insertion of exact duplicate copies
both preserve the set of contents and are covered by this reasoning.
This claim presupposes validity of both representations.
Duplicating an authenticated record in a way that invalidates
a signed envelope is outside that premise and triggers
snapshot-level validation instead.

\item \textbf{Isolation from evidence for other candidates.}
Let $E'$ be a valid enlarged snapshot with every newly added
receipt bound to a candidate $d\ne c$, and every added review
bound exclusively to such receipts.
Fix $q\in Q_p$.
The exact-candidate conjunct fails for each added receipt,
so $A_q$ and $B_q$ are unchanged.
Each review which qualifies for a receipt in $A_q$ must
name that exact receipt.
A review bound exclusively to a receipt for $d\ne c$
therefore cannot become a new member of $V_r$
for any current $r\in A_q$.
The sets $V_r$ are unchanged, hence so are
$A_q^+$, $P_q$, and $F_q$.
Applying this for each required $q$ leaves
$C$, $N$, and $M$ unchanged, which gives the first
isolation claim.
The hypothesis excludes the addition of a review that
endorses a current receipt. Such a review can supply the
missing witness for a previously unreviewed pass.

\item \textbf{Change of candidate.}
Let $d=(s',a',p')$ and suppose $d\ne c$.
Componentwise equality of tuples gives the negation of
the conjunction $s'=s$, $a'=a$, and $p'=p$.
At least one of these three equalities fails.
In particular, a one-component change suffices.
Under the new policy the check set may itself change,
so use $Q_{p'}$ in the gate for $d$.
If validation fails, the decision is \textsc{Blocked}.
Otherwise choose any $q\in Q_{p'}$, which is possible
because this set is nonempty.
A receipt for $c$ cannot be eligible at $(d,q)$
because its candidate field differs.
A review of such a receipt cannot supply an eligible pass
for $d$, because its endorsement preserves the exact
receipt identity.
If the snapshot contains only evidence for $c$,
then both $A_q$ and $B_q$ are empty at $d$.
It follows that $P_q$ and $F_q$ are both false.
The same argument applies to every $q\in Q_{p'}$.
Hence $C=N=\varnothing$, while $M=Q_{p'}\ne\varnothing$.
The third valid-input case yields \textsc{Blocked}.
This applies to artifact-only changes after a rebuild,
source-only changes even if the artifact is identical, and
policy-only changes with both source and artifact unchanged.
A repeated use of the very same $c$ is outside the change
hypothesis. A time-based expiry requirement, if desired,
must be specified and bound by the policy.

\item \textbf{Boundary cases and sharpness.}
The following cases exercise each boundary of the predicates
and show why the premises and precedence are useful.
\begin{enumerate}
\item \emph{Zero checks.}
With $Q_p=\varnothing$, all three witness sets would be empty
if formed. The decision is \textsc{Blocked} by validation,
so an empty required-check manifest cannot create eligibility.
\item \emph{Zero records.}
With $Q_p\ne\varnothing$ and $E=\varnothing$,
there are no eligible receipts and no reviews.
For every $q$, $(P_q,F_q)=(0,0)$.
Thus $M=Q_p$ and $C=N=\varnothing$.
The result is \textsc{Blocked} for missing evidence.
\item \emph{One check and one eligible failure.}
Let $Q_p=\{q\}$ with one eligible failure and no reviews.
Then $B_q$ is a singleton, $A_q^+$ is empty,
and the gate is \textsc{Rejected}.
This is the least nonempty failure witness case.
\item \emph{One check and producer self-review.}
Let $Q_p=\{q\}$ with an eligible pass $r$ and only its
producer's endorsement.
Distinctness fails, $V_r=\varnothing$, and $P_q$ is false.
With no failure, the decision is \textsc{Blocked}.
Adding a qualifying endorsement for the same $r$
makes $P_q$ true and yields \textsc{Ready}.
An endorsement of a different receipt leaves this case unchanged.
\item \emph{Minimum complete success.}
Let $Q_p=\{q\}$ with exactly one eligible pass and one
qualifying review of it and no failure.
Then $A_q^+$ is a singleton and the gate is \textsc{Ready}.
Deleting either required witness removes $P_q$;
an endorsement remaining after receipt removal cannot
supply the missing receipt.
\item \emph{Mixed required-check coverage.}
For $Q_p=\{q_1,q_2\}$, let only $q_1$ have a reviewed pass
and neither check have a failure.
Then $M=\{q_2\}$, so one success does not
satisfy the universal condition.
If an eligible failure for $q_1$ is then added,
$C=\{q_1\}$, still \textsc{Blocked}, now for conflict.
If the failure is instead added for $q_2$,
$N=\{q_2\}$ and $C=\varnothing$, giving \textsc{Rejected}.
\item \emph{Maximal presence of conflicting evidence.}
If every required check has both kinds of witnesses,
then $C=N=Q_p$ and $M=\varnothing$.
Since $Q_p\ne\varnothing$, the decision is
\textsc{Blocked}, despite complete reviewed-pass coverage.
At the other extreme one conflicting check suffices.
\item \emph{One material component changes.}
Keep two of $s$, $a$, and $p$ fixed and change the third.
Tuple inequality follows in each of the three possibilities.
Even if every check on the old candidate had a reviewed pass,
no such receipt witnesses a pass on the new candidate.
The transfer argument requires exact binding to all
three components.
\item \emph{Late authenticated failure.}
Let a valid snapshot be ready and add an eligible failure
for some required check $q$ while keeping its reviewed pass.
Then $q\in C$ in the enlarged snapshot.
The new decision is \textsc{Blocked}.
The earlier decision pertains to the earlier snapshot.
This case shows why readiness cannot be reused as a
prediction about later evidence.
\item \emph{Unlisted and malformed records.}
A receipt for a check outside $Q_p$ cannot enter
any required-check set.
A malformed receipt for a required check cannot become
eligible solely by naming the right check.
An unauthenticated snapshot blocks the entire gate even
when some contained bytes would describe passing results.
These three cases separate relevance, individual eligibility,
and snapshot validity.
\end{enumerate}

\item \textbf{Finite offline decision procedure.}
One can compute the predicates by finite scans.
Let $n$ be the number of record occurrences in $E$,
including duplicates, and let $k=|Q_p|$.
Both are nonnegative integers, with $k\ge 1$.
An occurrence is a position in a local representation of $E$.
For each $q$, use the following construction.
\begin{enumerate}
\item Set two Boolean flags to false.
The first represents the existence of an independently
reviewed pass found so far.
The second represents the existence of an eligible failure
found so far.
The zero-position prefix is empty, so both statements
are true descriptions of their respective flags.
\item Scan the $n$ record positions.
A position that is not an eligible receipt for $(c,q)$
leaves both flags unchanged.
An eligible failure sets the second flag and does not
change the first.
An eligible pass initiates a scan over the same finite
snapshot for a qualifying review of that pass.
\item Within the review scan, start with a false review flag.
For each position, test qualification against the currently
selected receipt.
Update the review flag by logical disjunction with that test.
For a prefix of length zero the flag is false by initialization.
If it correctly describes existence in a prefix of length $j$,
then after the next position it is true exactly when
there is a witness in that prefix or at the new position.
These are the only possibilities for a witness in the
prefix of length $j+1$.
Induction through length $n$ gives the qualifying-review
existence test on the entire snapshot.
\item For a pass, update the first main flag by disjunction
with the final review flag.
After any main prefix of length $i$, the first main flag is true
exactly when a pass in the prefix has a review somewhere
in the full snapshot.
The second main flag is true exactly when the prefix contains
an eligible failure.
The empty prefix gives the base case.
The three branches for the next record extend the invariant
by considering whether the next position is a new witness
of either class.
They exhaust receipt and result possibilities by eligibility,
which requires the result to be pass or fail.
Induction through length $n$ identifies the final flags
with $P_q$ and $F_q$.
\item After repeating across the $k$ required checks, form
$C$, $N$, and $M$ from the flags and apply the stated
precedence. These final sets can be recorded alongside
the decision, so a conflict does not erase evidence of
another required check being missing.
\end{enumerate}
There are at most $kn$ receipt eligibility tests during the
outer scans.
Each scanned position can trigger at most one length-$n$
review scan, so there are at most $kn^2$ qualification tests.
The combined count of these two types of tests satisfies
\begin{equation}
\label{eq:offline-ci-tests}
T_{\mathrm{scan}}\le k(n+n^2).
\end{equation}
For $n=0$ the right side is zero and there are no such tests,
even though initialization and the final decision still run.
For $n=1$ a snapshot with one pass receipt may initiate
one qualification test, but that sole receipt cannot qualify
as a separate review.
The count describes predicate evaluations after validation.
It is not a wall-time bound and omits manifest parsing,
signature verification cost, and the work of the checks.
The local scan requires all verification data to be provided
locally. If a verification dependency is unavailable,
validation cannot assert a ready result.

\item \textbf{Minimal information for this gate.}
The following indistinguishability arguments use the
principle that a deterministic evaluator has the same
result on equal inputs.
They apply to a simplified evaluator that deliberately drops
one of the specified distinctions.
Each comparison can be realized with a singleton required
set, an authenticated snapshot, and authorized identities.
\begin{enumerate}
\item Drop the artifact component from candidate binding.
Use a candidate $c$ and evidence that makes $c$ ready.
Evaluate the same evidence with an artifact-only change
to $c$ and unchanged source and policy.
The projected candidate is equal in the two evaluations,
and the evidence and policy inputs are unchanged.
The simplified evaluator therefore cannot distinguish them.
The exact gate requires \textsc{Ready} for the first and
\textsc{Blocked} for the second by candidate isolation.
\item Drop the source component instead and use a source-only
change, keeping artifact and policy fixed.
The projection is again equal and the input evidence is
unchanged.
The exact gate distinguishes the two decisions by the same
exact-binding test, even for builds producing equal outputs.
\item Drop the policy component from receipt binding.
Consider two policy manifests with the same required check,
authorized identities, and check configuration, but different
release-policy metadata declared material.
Fix the source, artifact, and original evidence.
A gate that has discarded all evidence about which policy
was used cannot use that receipt to tell whether it was
created under the original policy or the new policy.
More explicitly, keep the new candidate fixed and compare
two evidence collections that differ only in which policy is
named by the receipt and its bound review.
Under a projection that erases that distinction and keeps
only the remaining logical fields, their evidence is equal.
One collection makes the new candidate ready, and the
other lacks evidence for it.
This comparison concerns logical binding fields after
authentication. It does not assume authenticators or raw
record encodings remain identical.
\item Drop the eligible-failure bit while retaining the
reviewed-pass bit.
Compare a snapshot containing a reviewed pass with one
that adds an eligible failure.
The reviewed-pass bit stays true.
The exact result changes from \textsc{Ready} to
\textsc{Blocked}. Consequently the pass bit alone is
insufficient, even with only one required check.
\item Drop the reviewed-pass bit while retaining the
eligible-failure bit.
Compare a reviewed pass with a snapshot missing its review.
Both have false failure flags.
The results are respectively \textsc{Ready} and
\textsc{Blocked}, so the failure bit alone is insufficient.
\end{enumerate}
These examples show necessity of the distinctions for the
stated gate decision under the specified projections.
They do not claim minimal byte length for a receipt format,
or rule out an alternative encoding of the same information.

\item \textbf{Snapshot binding under change.}
With fixed $c$ and unchanged eligibility and review rules,
enlarging a valid snapshot can only preserve or add
witnesses for $P_q$ and $F_q$.
Indeed, membership of an old eligible receipt and of its
old qualifying review is preserved, so their conjunction is
preserved.
Taking an existential quantifier over a larger collection
cannot turn a true predicate of this kind false.
This gives information growth at the predicate level.
Readiness itself does not have that growth property.
The late-failure case proves failure of upward closure.
In the reverse direction, adding a qualifying review to an
eligible pass with no failure can change \textsc{Blocked}
to \textsc{Ready}, so readiness is not downward closed either.
A later authenticated snapshot may exclude records under
a revocation rule rather than merely append them.
The theorem applies separately to that new validated
snapshot. The witness-preservation claim does not apply
when eligibility changes or a record is removed.
For local use the evaluation and its evidence record must
both identify the snapshot instance.
An implementation must use a consistent snapshot at the
decision point, or validate the relevant identities again
before relying on the decision.
No assumption in the theorem makes a separately
executed read and release step atomic.

\item \textbf{Independence audit and verification obligations.}
The implications use finite-set membership,
Boolean logic, and exact equality.
They use no probabilistic independence, sample counts,
cryptographic injectivity, wall-clock ordering, or
network service status.
Distinct authenticated identities give a separation condition
on receipt authorship and review.
They do not provide evidence about reviewer judgment,
shared infrastructure, or collusion.
A local implementation seeking to apply the theorem must
first verify the following obligations on the evidence
it proposes to accept:
\begin{enumerate}
\item Confirm that the exact candidate manifests bind every
material input required by policy, including the material
toolchain identities when applicable.
\item Validate the snapshot as a whole, verify authorizations,
and check that the required-check set is nonempty.
\item Reconstruct the eligible receipts and verify each
review against its exact receipt and distinct producer.
\item Recompute the three sets in
\eqref{eq:offline-ci-conflict}--\eqref{eq:offline-ci-missing}
and reproduce the ordered decision.
\item Bind use of the result to the same candidate and
snapshot. If either changes, validate the intended pair
and evaluate again before use.
\end{enumerate}
These obligations are premises for applying the result,
not consequences of the decision.
In particular the theorem leaves the quality of individual
tests, operational attestation trust, and publication
authority to separately evidenced policy and human approval.
\qedhere
\end{enumerate}
\end{proof}

\begin{table}[!htbp]
\centering
\small
\caption{Offline evidence gate audit cases and their scope.}
\label{tab:offline-ci-audit}
\begin{tabularx}{\linewidth}{@{}p{0.28\linewidth}X p{0.15\linewidth}@{}}
\toprule
Case & Decision or limit & Source \\
\midrule
Empty required set & Blocked during validation. & \ref{def:offline-ci-snapshot} \\
Unreadable snapshot & Blocked during validation. & \ref{def:offline-ci-snapshot} \\
No eligible receipts & Blocked for missing pass evidence. & \ref{thm:offline-ci-acceptance} \\
Producer self-review & Cannot establish a reviewed pass. & \ref{eq:offline-ci-reviewed} \\
Review of another receipt & Cannot endorse this pass. & \ref{eq:offline-ci-reviewed} \\
Failure without review & Veto applies upon eligibility. & \ref{eq:offline-ci-fail} \\
Conflicting eligible results & Blocked even with complete pass coverage. & \ref{eq:offline-ci-conflict} \\
Missing check and failure & Rejected if there is no conflict. & \ref{thm:offline-ci-acceptance} \\
Complete reviewed passes & Ready if no eligible failure exists. & \ref{eq:offline-ci-ready} \\
Duplicate or permuted records & Same result for valid representations. & \ref{thm:offline-ci-acceptance} \\
Artifact-only change & Old candidate's records do not transfer. & \ref{thm:offline-ci-acceptance} \\
Late eligible failure & Reevaluate the enlarged snapshot. & \ref{thm:offline-ci-acceptance} \\
Local scan count & Bounds predicate evaluations after validation. & \ref{eq:offline-ci-tests} \\
\bottomrule
\end{tabularx}
\end{table}
\clearpage
